\begin{afterword}
\summaryhead

The post begins with the physical path of a spoken word, from a friend's throat through the listener's ear to the auditory cortex. The result is that the listener knows what the friend is looking at. Without knowledge of the mechanism, the author says, this would look like telepathy. Language is more complicated than it looks, but the brain hides the complexity.

The author then argues from design. An animal that heard ``tiger'' and reasoned about what its fellows meant by the word would be eaten, so there should be no step between recognizing a word and activating its concept. In the blegg network of earlier posts, this is a direct link from the unit that hears ``Blegg!'' to the central unit. From the inside, the label and the concept then feel almost the same, and the meaning feels like a property of the word. The author calls this a case of E.~T. Jaynes's Mind Projection Fallacy.

This feeling, the post concludes, is what keeps Albert and Barry arguing about what ``sound'' really means, as if it were a fact about the world.

\responsehead

The post's thesis is that words feel as if they carry their meanings, and that this feeling keeps disputes over definitions going. The first half of the thesis is old. In the first chapter of \textsc{The Meaning of Meaning} (1923), Ogden and Richards wrote that words ``mean'' nothing by themselves, ``as every one now knows,'' and that the belief that they did was once universal. The post names Jaynes's Mind Projection Fallacy as the general pattern, and Jaynes introduced the term while discussing how grammar disguises statements about perception. What the post adds is an explanation of why the old belief feels natural, and an application to the Standard Dispute of the day before.

The explanation is an argument from design: an organism that deliberated between hearing a word and activating its concept would be eaten. It is offered with ``you might visualize,'' and it is the same kind of reasoning as the network of ``How an Algorithm Feels From Inside,'' whose notes discuss it.

One claim goes beyond the design argument. The post says that the brain, ``on most occasions,'' does not separate word and meaning at all, ``only bothering to separate the two while learning a new language, perhaps.'' The claim is hedged, but a familiar experience separates them outside language learning: the word on the tip of the tongue, where a person has the meaning and cannot find the word. Researchers have produced this state on purpose since Brown and McNeill's study of 1966, by giving people definitions and asking for the words.

The application to the dispute fits the case as the post set it up. Albert and Barry feel that their argument is about a fact. The post shows why a question about a word can feel like one, and it ends with the test that exposes the confusion: state an experiment whose result depends on the dispute.

In short: an old observation, that words do not carry their meanings, explained by a design argument that is not evidence about brains, together with a claim about when people separate word from meaning that the tip-of-the-tongue state does not fit.
\end{afterword}
